% Transcription — fonds Alexandre Grothendieck, shelfmark 82, batch 2 (pages 21-40).
% Established from the facsimile archives/batches/82/batch-02.pdf, cut from a
% copy of 82.pdf fetched through the project's relay
% (https://grothendieck-relay.onrender.com/source/82.pdf); PDF page k =
% archivists' page 20+k, no cover sheet. Per the skill transcribe-grothendieck.
% The folder is seventy-eight pages in four batches (1-20, 21-40, 41-60,
% 61-78), transcribed in parallel by four passes. Batches 1, 3 and 4 existed
% when this file was finished and their notation was followed (E, E' in
% duality, phi, phi', f, f' = Lambda^2 f, Delta, rho_E, rho_E', rho~_E, the
% bases X, H, Y and X~, H~, Y~ of batch 1; GP(1) as \mathrm{GP}(1),
% \mathfrak{g}, \widetilde{\mathfrak{g}}, \mathcal{M}, his round T as
% \mathcal{T} of batch 3). Batch 1 was not edited.
% Pass: Opus 5.5 (claude-opus-5-5), 2026-09-26 — first pass, unchecked against the pages by a human.
%
% Dating and title. \dating is the folder's own catalogue date, copied
% verbatim; the folder belongs to the group « Géométrie et topologie
% combinatoire » (69-89), but since it has a date of its own the group's range
% is not added. \foldertitle is the catalogue title without its final full
% stop.
%
% Pages skipped: 35, a blank sheet (only show-through). Page 32 is a pink
% cover with nothing on it but « (21 p) » in pencil top right (hand not
% certain) and no title; by the cover rule it takes no \page marker, and it is
% described in the \note opening p. 33. Pages summarised: none. Pages
% transcribed from his hand: 21-31, 33-34, 36-40, in blue felt-tip; the top
% block of p. 33 is in black ink. Pages 33-34 are half-sheets written in
% landscape; p. 31 is one-third of a page, p. 38 two-thirds.
% Drawings: one on p. 33 (regular polygon inscribed in a circle, points
% zeta_i, vectors e_1, e_2, alpha), described in a \note, not redrawn.
% Diagrams set as tikzcd: p. 26 (the two exact sequences of the coupled
% forms), p. 28 (88), p. 29 (89) and (92).
%
% His pagination: none on the leaves. Numbered runs: the equation run of
% batch 1 ((1)-(72)) continues here from (73) on p. 22 to (93) on p. 29; on
% p. 22 a first (73)-(75) is struck and (73)-(74) restarted; the batch-1
% cover's « (29 p) » runs from p. 2 to about p. 31, where this run breaks off
% (« il nous faut pour ceci ouvrir une parenthèse »). After the cover p. 32
% (« (21 p) »), pp. 33-34 carry no numbers, and p. 36
% opens a new run (1)-(2) (the proof on the centraliser of u in GP(V) that
% batch 3 pp. 41-44 continue). The remark on p. 25 says it « devrait venir
% avec discussion après formule (42) », i.e. batch 1.
%
% Order of the leaves. P. 39 continues p. 37 (g h g^-1 = lambda h, lambda^2
% = 1, then g h^2 g^-1 = h^2). P. 38 does not continue p. 37: its opening
% (« alors un s.-groupe étale … sous-groupe ouvert de (Z/pZ)*_S, qui
% coïncide avec … {±1} fibre par fibre ») seems to continue the end of
% batch 3 p. 41 (N/N° ⊂ (Z/pZ)^×_S left hanging). This is an editorial
% conjecture, said in a \note; the pages are kept in archivists' order.
%
% What the batch is about. pp. 21-31, end of the run « Algèbres de Lie de
% rang 3 associées aux formes bil. symétriques de rang 3 »: remarks on the
% example sl(2) -> gp(1) when 2 is invertible (f iso) or zero (kernels,
% centres, derived algebras); Casimir forms (73)-(77), Cas_E'(x',y') =
% 2 phi'(x',y'), Cas_E(x,y) = -2 Delta phi(x,y), so Cas_E(x,x) = 8 q(x) for a
% special orthogonal (q, omega); how the brackets change when omega, phi are
% rescaled (78)-(81); characteristic 2: forms of sl(2) and gp(1) as pairs
% (V, extension), coupled pairs, symplectic affine planes, automorphism
% groups of dim 5 and 6 (pp. 26-27); then G = SO(q) = Aut(E, q, omega), its
% « simply connected cover » G~ with kernel mu_2 (82)-(88), Lie algebras
% g~, g (89)-(92), and the canonical isomorphisms E ≃ g~, E' ≃ g (93), of
% which E' -> g is proved injective fibrewise via (72) (p. 31). pp. 33-34
% (after the cover p. 32): minimal polynomials F_n of 1 + 2cos(2 pi/n)
% (F_3 ... F_6) and a rotation u of a regular n-gon in the basis e_1, e_2;
% then the chain Alg Spéc(M) ≃ P(g) ≃ P^v(g~) ≃ [P^v = Dr(P)] ≃ Div_2^+(X),
% identified with the smooth commutative subgroups of dimension 1 of G, and
% for such a T: Cent(T) = T, X - Y a T-torsor, N(T)/T ⊂ Aut(T) containing
% -id. pp. 36-40: V of rank 2, u in End(V) nowhere scalar; the commutant of
% u is O_S + O_S u, and its image in GP(V) is P^1 - Q_(Tr u, det u), étale
% iff the eigenvalues are distinct (p. 36); smooth commutative subgroups H
% of relative dimension 1 vs. rank-1 subbundles h of g (p. 37); is H the
% centraliser of h, the normaliser condition g u g^-1 = lambda u^(±1), and
% characteristic 2 counterexamples (gug^-1 = u + lambda, p. 40).
%
% Hardest pages: 26-27 (fast prose, a large struck block on 27), 34 (dense
% half-sheet with oblique marginal notes), 37 and 40 (bottom halves
% cancelled by crossing strokes, read only in part). Notation fixed for the
% batch: E, E' and \varphi, \varphi', \rho_E, \rho_{E'}, \tilde\rho_E as in
% batch 1; \mathfrak{g} = Lie(G), \widetilde{\mathfrak{g}} = Lie(G~); his
% gothic f for Lie(F) as \mathfrak{f}; his round T as \mathcal{T}; his
% underlined V on p. 27 as \mathbb{V}; \underline{\mathrm{Aut}},
% \underline{\mathrm{D\acute{e}r}}, \underline{\mathrm{End}} for his
% underlined functors.

\documentclass[11pt,a4paper]{article}
% Le préambule commun à toutes les transcriptions du fonds Grothendieck.
%
% Il tient en une idée : **l'incertitude se compose, elle ne se résout pas.**
% Une transcription qui lisse un mot douteux a détruit la seule chose pour
% laquelle on la faisait — un lecteur doit pouvoir savoir, à chaque endroit,
% ce qui était sur le feuillet et ce qui a été ajouté. Les six macros qui
% suivent sont donc l'appareil critique, et non de la décoration.
%
% Les mêmes commandes sont reconnues par `scripts/render.mjs`, qui produit la
% vue de lecture du site. Une macro ajoutée ici doit l'être aussi là-bas,
% faute de quoi le rendu échouera bruyamment — ce qui est voulu : un rendu
% silencieusement faux est pire qu'un rendu absent.

\ProvidesPackage{grothendieck}[2026/08/08 Transcriptions du fonds Grothendieck]

\usepackage{amsmath,amssymb,amsthm}
\usepackage{mathrsfs}
\usepackage{stmaryrd}
% Les diagrammes commutatifs sont partout dans ce fonds : c'est la langue
% même de la géométrie algébrique de Grothendieck, et une transcription qui
% les rendrait en prose ne transcrirait rien.
\usepackage{tikz-cd}
% Français par défaut — c'est la langue des feuillets — mais l'anglais est
% chargé aussi : la traduction et le résumé sont des documents anglais, et la
% typographie française (espaces avant les deux-points) y serait une faute.
% Ces documents basculent par \selectlanguage{english} en tête de corps.
\usepackage[english,french]{babel}
\usepackage{fontspec}
\usepackage{geometry}
\geometry{a4paper,margin=25mm,marginparwidth=28mm}
\usepackage{marginnote}
\usepackage{xcolor}
% ulem, not soul. `soul`'s \st and \ul are built on a character-by-character
% scan that XeLaTeX's font machinery breaks: the document fails with an
% undefined control sequence rather than degrading. ulem does the same two jobs
% and is Unicode-engine safe. `normalem` stops it hijacking \emph, which the
% transcriptions use for Grothendieck's own emphasis.
\usepackage[normalem]{ulem}
\usepackage{hyperref}
\hypersetup{colorlinks=true,linkcolor=black,urlcolor=black,pdfborder={0 0 0}}

% --- Métadonnées du lot -----------------------------------------------------
% Reprises telles quelles par le script de rendu, qui les affiche en tête de
% la vue de lecture. Elles fixent ce que le fichier prétend couvrir : sans
% elles, une transcription flotte sans point d'attache dans un fonds de seize
% mille feuillets.

\newcommand{\folder}[1]{\gdef\grothFolder{#1}}
\newcommand{\batch}[1]{\gdef\grothBatch{#1}}
\newcommand{\leaves}[2]{\gdef\grothFirst{#1}\gdef\grothLast{#2}}
\newcommand{\foldertitle}[1]{\gdef\grothFoldertitle{#1}}
\gdef\grothFolder{?}\gdef\grothBatch{?}\gdef\grothFirst{?}\gdef\grothLast{?}\gdef\grothFoldertitle{}

% --- Le résumé --------------------------------------------------------------
%
% Il ouvre la lecture modernisée, sous le titre « Résumé », et il est composé
% en petit. Ce n'est pas un ornement : c'est le seul passage du document qui
% ne soit pas des mathématiques, et un lecteur doit pouvoir le reconnaître
% avant de l'avoir lu — puis le sauter s'il connaît déjà le sujet, ou s'y
% arrêter s'il ne le connaît pas. Le filet qui le suit marque la reprise du
% corps.

\newenvironment{resume}
  {\small\setlength{\parskip}{0.45em}}
  {\par\medskip\hrule\medskip}

% --- Mots-clés ---------------------------------------------------------------
%
% En anglais, en clôture du résumé de la lecture modernisée : le vocabulaire
% moderne sous lequel chercher ce que les feuillets construisent. Le manifeste
% du site (`npm run manifest`) les extrait pour étiqueter la cote entière —
% cette ligne est la source unique des tags, il n'y a pas de fichier à côté.
\newcommand{\keywords}[1]{\par\smallskip\noindent
  {\footnotesize\textsc{Keywords} --- \textit{#1}}\par}

% --- L'appareil critique ----------------------------------------------------

% Le numéro de feuillet, en marge. C'est l'ancre qui permet de régler un
% désaccord sur une formule en regardant *un* feuillet plutôt qu'une chemise
% de six cents. Le site s'en sert aussi pour tourner les pages du fac-similé
% quand on fait défiler la transcription.
\newcommand{\leaf}[1]{%
  \marginnote{\footnotesize\color{gray!70}\textsf{#1}}%
  \label{leaf:#1}\ignorespaces}

% Illisible : on ne devine pas. Un mot inventé qui se lit comme les autres est
% le pire résultat possible.
\newcommand{\ill}{\textcolor{red!60!black}{[\,\ldots\,]}}

% Lecture douteuse : le mot est donné, et signalé comme incertain.
\newcommand{\uncertain}[1]{\uline{#1}}

% Ajout de l'éditeur — une lettre manquante, un mot sous-entendu.
\newcommand{\add}[1]{\textcolor{blue!50!black}{[#1]}}

% Biffé par Grothendieck lui-même. Ce qu'il a rayé fait partie du document :
% une rature dit souvent où la pensée a changé de direction.
\newcommand{\struck}[1]{\sout{#1}}

% Note du transcripteur, sur l'état matériel du feuillet.
\newcommand{\note}[1]{\par\smallskip{\small\color{gray!60!black}\textit{[#1]}}\par\smallskip}

% Note marginale de Grothendieck, distincte du corps du texte.
\newcommand{\marginal}[1]{\par\smallskip\noindent\hspace{1em}%
  {\small\color{gray!60!black}#1}\par\smallskip}

% --- Titre ------------------------------------------------------------------

\AtBeginDocument{%
  \begin{center}
    {\large\bfseries Fonds Alexandre Grothendieck --- cote n\textsuperscript{o}~\grothFolder}\\[2pt]
    {\itshape\grothFoldertitle}\\[4pt]
    {\small feuillets \grothFirst--\grothLast\ (lot \grothBatch)}\\[2pt]
    {\footnotesize\color{gray!60!black}Transcription établie d'après le fac-similé numérisé
     par l'Université de Montpellier.\\
     Les lectures incertaines sont soulignées, les illisibles notées [\,\ldots\,],
     les ajouts entre crochets.}
  \end{center}
  \vspace{1em}}

\endinput


\folder{82}
\batch{2}
\pages{21}{40}
\dating{[vers 1982]}
\watermark{Édition de démonstration}
\foldertitle{SO(3) ≃ GP(1) : notes manuscrites (s.d.)}

\begin{document}

\page{21}
\note{la page poursuit l'exemple de la page 20 (batch 1) : $E \simeq \mathfrak{sl}(2)$, $E' \simeq \mathfrak{gp}(1)$, $f$, $f'$, $\rho_E$ ; elle appartient au cahier annoncé par la couverture de la page 1.}
\emph{Remarque.} Si 2 inv.\ sur $S$, alors $f$ est un isomorphisme d'algèbres
de Lie, \uncertain{$f$} permettant d'identifier $E$ et $E'$ (c'est aussi
l'identification habituelle d'une \ill{} quadratique et de son dual), et
$\rho_E$ est identifié à $\rho_{E'}$ ou à $\tilde{\rho}_E$ ou
$\tilde{\rho}_{E'}$ \ldots\ ; $f'$ (qui s'identifie à $2\,\mathrm{id}_E$) n'a
donc guère d'intérêt ! (Ces constatations s'appliquent plus généralement
chaque fois que $\Delta$ est inv.)
\note{avant « permettant », un symbole surchargé, lu $f$ avec doute (on peut y voir $\rho_E$ ou $\varphi$).}

Si par contre 2 est nul sur $S$, alors $f$ a comme noyau
$\mathcal{O}_S\cdot\tilde{H}$, qui est donc le \emph{centre} de
$E = \mathfrak{sl}(2)$. Quant à $f'$, il a comme noyau
$\mathcal{O}_S X \oplus \mathcal{O}_S Y$, qui est \add{alors} aussi
l'Algèbre \emph{dérivée} de $E' = \mathfrak{gp}(1)$, \struck{\ill{}} \add{qui}
est aussi l'image de $f$, alors que l'image de $f'$ est le noyau
$\mathcal{O}_S\tilde{H}$ de $f$. [Par contre le centre de $\mathfrak{gp}(1)$
est nul -- c'est vrai sur toute base $S$, comme on vérifie facilement --
tandis que le groupe \uncertain{dérivé} de $\mathfrak{sl}(2)$ est égal au
centre $\mathcal{O}_S\tilde{H}$] \uncertain{Ici} l'existence de $f'$ (cf.\
plus bas un contexte \uncertain{général}) \uncertain{pour} avoir un
intérêt\ldots\ Il faudrait par la suite expliciter la structure des formes
de $\widetilde{\mathfrak{g}} = \mathfrak{sl}(2)$ et
$\mathfrak{g} = \mathfrak{gp}(1)$ en car.~2, et la conclusion
\uncertain{qu'elles} sont trop triviales pour \ill{}
\note{« qui » est écrit au-dessus d'un mot biffé ; « groupe dérivé » est la lecture de l'encre, là où l'on attend l'algèbre dérivée. La dernière ligne, serrée contre le bas de la feuille, se lit mal.}

\marginal{\emph{NB} \uncertain{Également}, sa connaissance \ill{} permet
\uncertain{les} \ill{} \uncertain{de} \ill{} qui \ill{} \ill{} \ill{}
multiple \ill{} \uncertain{complexe} \ill{}}
\note{note écrite en oblique dans la marge gauche, en face des lignes sur $f'$ ; presque entièrement illisible.}

\page{22}
un intérêt (contrairement aux formes de $\mathrm{SL}(2)_S$ ou
$\mathrm{GP}(1)_S$, qui sont équivalentes et équivalent aux formes de
$\mathbb{P}^1_S$ \ldots).

Revenant au cas général, il peut être intéressant de calculer les formes
\struck{quadratiques} \add{\uncertain{de}} Casimir de $\mathfrak{g} = E'$ et
$\widetilde{\mathfrak{g}} = E$, définies par

\struck{(73) $\mathrm{Cas}_{E}(x, y) = \rho_{E'}$, $\mathrm{Tr}\,\rho_E$}
\struck{(73)}
\[
  \begin{cases}
    \mathrm{Cas}_{E'}(x') = \mathrm{Tr}\, \rho_{E'}(x')^2 \\
    \mathrm{Cas}_{E}(x) = \mathrm{Tr}\, \tilde{\rho}_{E}(x)^2
  \end{cases}
\]
(les formes bil.\ associées sont respectivement
$2\,\mathrm{Tr}\, \rho_{E'}(x')\rho_{E'}(y')$ et
$2\,\mathrm{Tr}\, \tilde{\rho}_E(x)\tilde{\rho}_E(y)$).
Le calcul donne :
\struck{(74)}
\[
  \mathrm{Cas}_{E'}(x') = \Delta\, \varphi(x, x)
\]
d'où, en remplaçant $(E, \omega, \varphi)$ par $(E', \omega', \varphi')$
\struck{(75)} $\mathrm{Cas}_{E}(x) = \ldots$
\note{tout ce qui précède, depuis la première ligne (73), est annulé par quatre grands traits obliques et par un trait qui contourne le bloc ; une première ligne (73) est en outre biffée à part, et la ligne (75) s'arrête sur le signe égal. Dans la première (73), l'indice de $\mathrm{Cas}$ est surchargé ($E$ sur $E'$). Au-dessus de $\varphi$ dans (74) un mot est noirci.}
\[
  (73) \qquad
  \begin{cases}
    \mathrm{Cas}_{E'}(x', y') = \mathrm{Tr}\bigl(\rho_{E'}(x') \circ \rho_{E'}(y')\bigr) \\
    \mathrm{Cas}_{E}(x, y) = \mathrm{Tr}\bigl(\tilde{\rho}_{E}(x) \circ \tilde{\rho}_{E}(y)\bigr)
  \end{cases}
\]
On trouve alors
\[
  (74) \qquad
  \begin{cases}
    \mathrm{Cas}_{E'}(x', y') = 2\varphi'(x', y') \\
    \mathrm{Cas}_{E}(x, y) = -2\Delta\, \varphi(x, y)
  \end{cases}
\]
\note{le chiffre des unités du numéro (74) est surchargé ; le signe $-$ de la seconde ligne est écrit par-dessus un premier signe.}
Pour le voir, on note d'abord que les \add{\ill{}} formules \ill{}, et la
propriété d'homogénéité qu'il faut, quand on multiplie $\omega$ par
$\lambda$, $\varphi$ par $\mu$ : les
\note{les deux dernières lignes, en bas de page, se lisent mal ; un mot illisible est inséré au-dessus de « formules ».}

\page{23}
\emph{deux} membres de la $1^{\mathrm{re}}$ sont multipliés tous par
$(\mu\lambda)^2 = \mu^2\lambda^2$, et les \emph{deux} membres de la seconde
par $(\mu^2\lambda)^2 = \mu^4\lambda^2$. On vérifie ensuite les formules
dans le cas particulier (21). On peut se simplifier le travail en
\uncertain{notant} que pour l'une et l'autre formule, les deux membres sont
des formes invariantes sous l'action des groupes \add{spéciaux} orthogonaux
$\underline{\mathrm{O}}(q, \omega)$ (sur $\operatorname{Spec} \mathbf{Z}$
\uncertain{si on veut}) or il n'y a (\uncertain{à} un \uncertain{facteur}
\uncertain{près}, \ill{} -- cela suffit -- qu'une seule forme invariante
sym., à scalaire près, ce qui montre la proportionnalité. Pour avoir le
coefficient correctif, il suffit de calculer la valeur des 2 membres pour un
couple d'arguments pour lesquels ils soient non nuls, p.ex.\ $x' = y' = H$,
$x = y = \tilde{H}$. On trouve
\[
  \begin{aligned}
    &\mathrm{Cas}_{E'}(H, H) = 1 + 1 + 0 = 2 \\
    &\varphi'(H, H) = 1 \\
    &\mathrm{Cas}_{E}(\tilde{H}, \tilde{H}) = \mathrm{Cas}_{E'}(2H, 2H) = 4\, \mathrm{Cas}_{E'}(H, H) = 8 \\
    &\qquad (\ldots = 4 + 4 + 0) \\
    &\varphi(\tilde{H}, \tilde{H}) = 2 \\
    &\qquad \Delta = -2
  \end{aligned}
\]
\note{la ligne « (sur Spec $\mathbf{Z}$ \ldots{} qu'une seule forme invariante sym., à scalaire près » se lit mal entre les formules ; la parenthèse ouverte avant « à un facteur près » n'est pas refermée. Devant « $= 4 + 4 + 0$ », deux mots lus avec doute « ou encore ».}
On note qu'en vertu du \uncertain{bref} principe de conservation des
identités polynômiales, il suffit de prouver (74) sur un corps de car.\
0 alg.\ clos, \struck{avec $f$ un iso i.e.\ $\varphi$ \ill{} \ldots}

\page{24}
à une forme quadratique non dég.\ -- quitte (grâce aux \uncertain{Cas}), pour
une base convenable, \uncertain{multiplier} par \uncertain{$\lambda$} la forme
$xy + z^2$ (quitte à multiplier $\omega$ par un scalaire \uncertain{convenable}
\uncertain{inversible}) on peut supposer que $\omega = e_1 \wedge e_2 \wedge e_3$,
ce qui nous ramène au cas déjà traité.
\note{les quatre premières lignes sont rapides et se lisent mal ; seules les formules sont sûres.}

Lorsque $\varphi$ est associé à une forme quadratique $q$, on tire de la
deuxième formule (74), \uncertain{compte} tenu de $\varphi(x, x) = 2q(x)$ :
\[
  (75) \qquad \mathrm{Cas}_E(x, x) = -2\Delta\, \varphi(x, x) = -4\Delta\, q(x)
\]
ce qui, dans le cas d'une forme \add{\uncertain{q}} spéciale orthogonale
$(q, \omega)$, \uncertain{où} $\Delta = -2$, donne
\[
  (76) \qquad \mathrm{Cas}_E(x, x) = 8q(x)
\]
Donc dans ce cas, \struck{\ill{}} \uncertain{si} 2 est \uncertain{inv.}
(\uncertain{ou} plus généralement non div.\ de 0) $q(x)$ se \emph{déduit} de
la structure d'Algèbre de Lie sur $E$, comme l'opér.\ de Casimir divisé par
8\ldots

On trouve de même, par la première formule
\[
  (77) \qquad \mathrm{Cas}_{E'}(x', x') = 2\varphi'(x', x') = 2q'(x')
\]
i.e.\ on retrouve la forme quadratique $q'(x')$ (ou la forme bil.\ $\varphi'$
dont elle provient) à l'aide de la structure d'Algèbre de Lie de
$E' = \mathfrak{g}$.
\note{le numéro (77) est écrit au-dessus d'un numéro noirci ; dans la dernière ligne, $E'$ est surchargé.}

\page{25}
\emph{Remarques} (devrait venir avec discussion après formule (42)). Le
crochet $[\ ,\ ]_{\varphi}$ \add{$= [\ ,\ ]_{\varphi, \omega}$} sur $E'$ est
multiplié par $\lambda\mu$ si $\omega$ est multiplié par $\lambda$, $\varphi$
par $\mu$ ; \struck{\ill{}}
\note{« $= [\ ,\ ]_{\varphi, \omega}$ » est inséré au-dessus de la ligne, devant « sur $E'$ ».}
\[
  (78) \qquad
  \begin{cases}
    [\ ,\ ]_{\lambda\omega, \mu\varphi} = \lambda\mu\, [\ ,\ ]_{\omega, \varphi} \\
    [\ ,\ ]_{\lambda\omega, (\mu\varphi)'} = \lambda\mu^2\, [\ ,\ ]_{\omega, \varphi'}
  \end{cases}
\]
On en conclut que
\[
  (79) \qquad [\ ,\ ]_{\omega, \varphi} = [\ ,\ ]_{\omega_1, \varphi_1}
  \iff \exists \lambda \in \Gamma(S, \mathcal{O}_S^{*}),\ \omega_1 = \lambda\omega,\ \varphi_1 = \lambda^{-1}\varphi
\]
\note{la condition à droite de (79) est écrite dans un encadré ouvert ; le $\Gamma(\ldots)$ au-dessus porte un symbole mal formé, lu $\mathcal{O}_S^{*}$ par le sens.}
\struck{D'autre part}
(compte tenu de ce que pour $\omega$ fixé, la corr.\ entre $\varphi$ et
$[\ ,\ ]_{\varphi}$ est biunivoque). D'autre \add{part} \add{et pour $\omega$
fixé}, si $\det(f) = \delta(\varphi)$ \emph{inversible} i.e.\ $f$ iso, alors
$f \mapsto f' = \alpha'_{\omega} \circ \Lambda^2 f \circ \alpha_{\omega}^{-1}$
est \struck{bijectif} \add{\uncertain{presque}} \uncertain{injectif},
\struck{\ill{}} plus précisément si $f, g : E \to E'$, on a
\[
  (80) \qquad f' = g' \iff g = \varepsilon f , \quad \varepsilon \in \Gamma(S, \mu_2) \ \text{i.e.}\ \varepsilon^2 = 1
\]
\note{la ligne « D'autre part et pour $\omega$ fixé » est une insertion interlinéaire ; au-dessus de « bijectif » biffé, un mot lu avec doute « presque ».}
[car $f' = g' \Rightarrow \Delta(f) = \Delta(g)$ i.e.\ $\Delta(f)^2 = \Delta(g)^2$
i.e.\ $\Delta(g) = \varepsilon\Delta(f)$, d'où compte tenu de
$f = \Delta(f)^{-1}(f')'$, $g = \Delta(g)^{-1}(g')'$ \ldots\ $g = \varepsilon f$ ]
On en déduit aisément que
\[
  (81) \qquad [\ ,\ ]_{\omega_1, (\varphi_1)'} = [\ ,\ ]_{\omega, (\varphi)'}
  \iff \omega_1 = \lambda\omega ,\ \varphi_1 = \mu\varphi
\]
avec $\lambda\mu^2 = 1$ i.e.\ $\mu$ \emph{une} racine carrée de $\lambda^{-1}$
(\emph{NB} $\lambda$ et $\mu$ uniquement déterminés).

\page{26}
En car.~2, la catégorie des formes de $\mathfrak{sl}(2)$ est équivalente à
celle des couples $(V, E)$ où $V$ est un Module \uncertain{vectoriel} de
rang 2 et $E$ une extension de \struck{\ill{}} $V$ par $\det V$
(\emph{NB} pour $V$ fixé, la donnée d'une telle extension équivaut à celle
d'une \struck{\ill{}} extension de $\det V^{-1}$ par $\check{V}$ ou de
$\mathcal{O}_S$ par $(\det V) \otimes \check{V} \simeq V$, i.e.\ à la cat.\
des fibrés affines sur $V$\ldots). La catégorie des formes de
$\mathfrak{gp}(1)$ équivaut à celle des \struck{\ill{}} couples $(V', E')$
d'un Module \uncertain{vectoriel} de rang 2, et d'une extension $E'$ de
$\mathcal{O}_S$ par $V'$ ; [pour $V'$ fixé, la cat.\ équivaut à celle des
fibrés affines sur $V'$).
\note{au-dessus de $V'$ dans « couples $(V', E')$ », un $\check{V}'$ est écrit en interligne, qui peut corriger $V'$.}

Disons que deux telles algèbres de Lie $E, E'$ sont « couplées » si on a un
accouplement dualisant entre $E$, $E'$, dans lequel le sous-module $\det V$
de $E$ est l'orthogonal du sous-module $V'$ de $E'$

\begin{tikzcd}[column sep=small]
    0 \arrow[r] & \det V \arrow[r] & E \arrow[r] & V \arrow[r] & 0 \\
    0 & \mathcal{O}_S \arrow[l] & E' \arrow[l] & \check{V}' \arrow[l] & 0 \arrow[l]
\end{tikzcd}

\note{au-dessus de $\det V$, relié par $\wr$, est écrit $\mathcal{O}_S$ ; sous $\check{V}'$, relié par $\wr$, est écrit $V$. Le $\mathcal{O}_S$ de la seconde ligne est surchargé.}
-- $V'$ est donc \struck{\ill{}} $\check{V}$ -- et donc le dual de
$\mathcal{O}_S$, dual de $V$, est $\det V^{-1}$ \uncertain{d'où} \ill{} de
$\mathcal{O}_S$. Donc $V$ est \add{\uncertain{à équiv.}} \uncertain{fait}
muni d'un \uncertain{iso} $\det V \simeq \mathcal{O}_S$ \struck{\ill{}} d'une
structure symplectique, et $V' \simeq V$. La catégorie des couples
\add{accouplés} \struck{de} formes de $\mathfrak{sl}(2)$, $\mathfrak{gp}(1)$
est donc équivalente à celle des triples $(V, \varphi, E)$ d'un Module
\uncertain{vectoriel} $V$ de rang 2, une structure symplectique $\varphi$ et
ext.\ de $V$ par $\mathcal{O}_S$ (ce qui revient à ext.\ de $\mathcal{O}_S$
par $\check{V} \simeq V$, ou encore \uncertain{à un} fibré affine sur $V$) :
c'est donc équivalente
\note{les lignes qui suivent le diagramme sont très corrigées (insertions « à équiv. », « accouplés ») et leur prose se lit mal ; « Module vectoriel » est la lecture de la page 36, où il écrit « Module (fibré) vectoriel ».}

\page{27}
à la donnée d'un fibré en « plans affines symplectiques » (en
\uncertain{car.\ deux}). Le \add{\uncertain{schéma en}} groupes des
automorphismes d'une telle structure est donc le produit semi-direct de
$\underline{\mathrm{SL}}(V)$ et de $\mathbb{V}$ -- lisse de dim~5 -- \ill{}
\uncertain{isomorphe} \ill{} \ill{} forme de $\mathrm{GP}(1)$ ! Le groupe
\uncertain{des} \ill{} forme de $\mathfrak{sl}(2)$ \uncertain{ou} \ill{} des
automorphismes d'une forme de $\mathfrak{gp}(1)$ est $\simeq$ \uncertain{groupe}
\uncertain{affine} \struck{\ill{}} \ill{}, lisse de dim rel.~6.
\note{la fin du premier paragraphe se lit mal ; au-dessus de « groupe » est inséré un mot lu avec doute « schéma en ». « $\mathbb{V}$ » est un $V$ souligné.}

On déduit aisément, de la structure précédente, les \uncertain{applications}
$f$, $f'$, $\rho_{E'}$. En effet, $f$ est le composé
$E \xrightarrow{\mathrm{can}} V \xrightarrow{\varphi} V' \xrightarrow{\mathrm{can}} E'$,
$f'$ est le composé
$E' \xrightarrow{\mathrm{can}} \mathcal{O}_S \xrightarrow{\varphi} \det V \xrightarrow{\mathrm{can}} E$,
\struck{\uncertain{enfin} $\rho_{E'}$ (\uncertain{est} l'opération \uncertain{qui}
\ill{} $V' = f(E)$ \ill{} de $E'$ \ldots\ \ill{} la représentation adjointe de
$E$, qui est nulle sur $\det V = \mathcal{O}_S \subset E$, et qui \ill{} par un
$H \in E'$ au-dessus de 1 donne l'opération $\rho_{E'}(H) : E \to E$ qui
(\uncertain{étant} nulle sur $\det V = \mathcal{O}_S$) se factorise en
\ill{} \uncertain{symplectique})}
\note{à partir de « enfin $\rho_{E'}$ », tout le bas de la page est annulé par un trait qui l'encadre et deux grands traits obliques ; au-dessus de la ligne, « (op.\ de $E'$ sur $E$) » est inséré après $\rho_{E'}$. Le reste de la feuille est blanc.}

\page{28}
Soit $q$ forme quadratique lisse sur $E$, $\varphi$ la forme bilinéaire
associée, $\omega$ une base de $\underline{\det} E$ telle que
\[
  (82) \qquad \delta'(q) = -\omega^{\otimes(-2)}
\]
Posons
\[
  (83) \qquad G = \underline{\mathrm{SO}}(q) = \underline{\mathrm{Aut}}(E, q, \omega)
\]
Soit $\widetilde{G}$ le \struck{\ill{}} « revêtement \uncertain{simplement}
connexe » de $G$, \struck{\ill{} canoniquement} i.e.\ le noyau de
$\widetilde{G} \to G$, canoniquement isomorphe à $\mu_2$ \struck{\ill{} des
racines de \uncertain{l'unité}}, posons
\[
  (84) \qquad 1 \to \mu_2 \to \widetilde{G} \xrightarrow{F} G \to 1
\]
Soient
\[
  (85) \qquad
  \begin{cases}
    \widetilde{\mathfrak{g}} = \underline{\mathrm{Lie}}(\widetilde{G}) \\
    \mathfrak{g} = \underline{\mathrm{Lie}}(G)
  \end{cases}
\]
Alors on a un hom.\ can.\ (épim.\ de noyau $\mu_2$)
\[
  (86) \qquad F : \widetilde{G} \to G
\]
et une opération obtenue par passage au quotient dans la représentation
adjointe de $\widetilde{G}$ sur lui-même
\[
  (87) \qquad G \xrightarrow{\ \rho\ } \underline{\mathrm{Aut}}_{S\text{-}\mathrm{gr}}(\widetilde{G})
\]
rendant commutatif le diagramme (88)

\begin{tikzcd}
  \widetilde{G} \arrow[r, "F"] \arrow[dr, "\mathrm{ad}_{\widetilde{G}}"'] & G \arrow[d, "\rho_{\widetilde{G}}"] \arrow[dr, "\mathrm{ad}_G"] & \\
  & \underline{\mathrm{Aut}}_{S\text{-}\mathrm{gr}}(\widetilde{G}) \arrow[r] & \underline{\mathrm{Aut}}_{S\text{-}\mathrm{gr}}(G)
\end{tikzcd}

\note{la flèche horizontale du bas porte « pass.\ au quotient ». Dans (84), (87) et (88), $G$ porte en exposant un petit signe surchargé, non interprété ; dans (84), après $G$, un second signe est biffé. Le noyau de $\widetilde{G} \to G$ est dit « canoniquement isomorphe à $\mu_2$ » en interligne, au-dessus d'une ligne biffée.}

\page{29}
ou ce qui revient au même, on a comm.\ dans
\note{(89) : diagramme à trois étages, numéroté (89).}

\begin{tikzcd}
  \widetilde{G} \times \widetilde{G} \arrow[r, "\mathrm{ad}_{\widetilde{G}}"] \arrow[d, "F \times \mathrm{id}_{\widetilde{G}}"] & \widetilde{G} \arrow[d, "\mathrm{id}"] \\
  G \times \widetilde{G} \arrow[r, "\rho"] \arrow[d, "\mathrm{id}_G \times F"] & \widetilde{G} \arrow[d, "F"] \\
  G \times G \arrow[r, "\mathrm{ad}_G"] & G
\end{tikzcd}

\note{les $G$ de la colonne de gauche et de la dernière ligne portent le même petit exposant surchargé qu'à la page 28 ; la flèche de droite du haut est doublée d'un signe mal formé, lu $\mathrm{id}$ avec doute.}

En passant aux algèbres de Lie, on trouve
\[
  (90) \qquad \widetilde{\mathfrak{g}} \xrightarrow{\ \mathfrak{f}\ } \mathfrak{g} \qquad \text{hom.\ d'alg.\ de Lie}
\]
et une opération
\[
  (91) \qquad \mathfrak{g} \xrightarrow{\ \rho_{\mathfrak{g}}\ } \underline{\mathrm{D\acute{e}r}}(\widetilde{\mathfrak{g}})
\]
\note{la flèche de (90) porte une lettre gothique surchargée, lue $\mathfrak{f}$ ; celle de (91) une lettre $\rho$ d'indice mal formé.}
rendant commutatif le diagramme (92)

\begin{tikzcd}
  \widetilde{\mathfrak{g}} \times \widetilde{\mathfrak{g}} \arrow[r, "\mathrm{ad}_{\widetilde{\mathfrak{g}}}"] \arrow[d, "\mathfrak{f} \times \widetilde{\mathfrak{g}}"] & \widetilde{\mathfrak{g}} \arrow[d, "\mathrm{id}"] \\
  \mathfrak{g} \times \widetilde{\mathfrak{g}} \arrow[r] \arrow[d] & \widetilde{\mathfrak{g}} \arrow[d, "\mathfrak{f}_{\widetilde{\mathfrak{g}}}"] \\
  \mathfrak{g} \times \mathfrak{g} \arrow[r, "\mathrm{ad}_{\mathfrak{g}}"] & \mathfrak{g}
\end{tikzcd}

\note{devant la première ligne, un signe biffé. L'indice de la flèche de droite du bas est lu avec doute.}

\struck{Remarques} Ceci posé, on va définir des isom.\ canoniques de Modules :
\[
  (93) \qquad
  \begin{cases}
    E \xrightarrow[\sim]{\ i\ } \widetilde{\mathfrak{g}} \\
    E' \xrightarrow[\sim]{\ j\ } \mathfrak{g}
  \end{cases}
\]
\struck{\ill{}} Transportant par ces isom.\ les lois $[\ ,\ ]$ sur
$\widetilde{\mathfrak{g}}$ et $\mathfrak{g}$, et
$\mathfrak{f} : \widetilde{\mathfrak{g}} \to \mathfrak{g}$
\note{la seconde flèche de (93) porte une lettre lue $j$, comme à la page 30.}

\page{30}
et $\rho_{\mathfrak{g}} : \mathfrak{g} \to \underline{\mathrm{D\acute{e}r}}(\widetilde{\mathfrak{g}})$,
on trouve respectivement : la loi $[\ ,\ ]_{\varphi'}$ sur $E$, la loi
$[\ ,\ ]_{\varphi}$ sur $E'$, l'application $f$ \add{$(E \to E')$}, et l'hom.\
$\rho_E$.
\note{« $(E \to E')$ » est écrit au-dessus de la ligne, relié à $f$ par un trait.}
Enfin, l'hom.\ composé
\[
  \mathfrak{g} \xrightarrow{\ \rho_{\mathfrak{g}}\ } \underline{\mathrm{D\acute{e}r}}(\widetilde{\mathfrak{g}})
  \xrightarrow[(i)]{\ \sim\ } \underline{\mathrm{D\acute{e}r}}(E) \to \underline{\mathrm{End}}(E)
\]
n'est autre que la repr.\ linéaire de $\mathfrak{g} = \underline{\mathrm{Lie}}(G)$
correspondant à la repr.\ \struck{\uncertain{tautologique}} \uncertain{linéaire}
de $G$ sur $E$.
\note{au-dessus de la première flèche, un mot noirci ; « tautologique » est biffé et le mot qui le remplace se lit mal.}

\emph{Construction.} Considérons $E'$ avec la loi $[\ ,\ ]_{\varphi}$, on sait que
$E'$ opère sur $E$ par $\rho_{E'}$
\[
  E' \longrightarrow \underline{\mathrm{D\acute{e}r}}(E, [\ ,\ ]_{\varphi'})
\]
en fait on sait que $E'$ s'envoie dans les endom.\ \add{de $E$} qui invariant
infinitésimalement $q$ \add{(et qui invariant $\omega$)}, lesquels forment
précisément l'algèbre de Lie $\mathfrak{g}$ de $G$, d'où finalement un homom.
\[
  E' \xrightarrow{\ j\ } \mathfrak{g}
\]
qui est compatible avec le crochet. Je dis \emph{que c'est un iso}. Comme
$E'$ et $\mathfrak{g}$ sont \struck{\ill{}} loc.\ libres de rang 3, il suffit
de prouver que c'est injectif sur chaque fibre, \add{\uncertain{point}}\ldots
\note{« (et qui invariant $\omega$) » est écrit au-dessus de la ligne, entre parenthèses ; au-dessus de « fibre », un mot inséré lu avec doute.}

\page{31}
encore que $E'$ opère \emph{fidèlement} sur $E$ par $\rho_{E'}$, quand
$S = \operatorname{Spec} k$, $k$ corps alg.\ clos. Or dans ce cas on peut
trouver une base $e_1, e_2, e_3$ p.r.\ à laquelle $q$ se met sous la forme
$xy + z^2$, $\omega = e_1 \wedge e_2 \wedge e_3$ (quitte \add{au besoin} à
changer $e_3$ en $-e_3$\ldots) et il suffit de prendre alors les (\uncertain{calculs}
de (72))
\[
  \rho_E(xX + hH + yY) :
  \begin{cases}
    \tilde{X} \mapsto h\tilde{X} - y\tilde{H} \\
    \tilde{Y} \mapsto x\tilde{H} - h\tilde{Y} \\
    \tilde{H} \mapsto -2x\tilde{X} + 2y\tilde{Y}
  \end{cases}
\]
\[
  \rho_E(xX + hH + yY) = 0 \Longrightarrow x = y = h = 0 \qquad \text{o.k.}
\]
\note{dans la deuxième ligne de l'accolade, le signe devant $h\tilde{Y}$ est surchargé ($+$ ou $-$) ; la formule (72) de la page 20 donne $-$.}
Il faut encore définir l'isom.
\[
  E \xrightarrow{\ \sim\ } \widetilde{\mathfrak{g}}
\]
et prouver ses propriétés. Pour ceci il faut redéfinir avec précision, en
termes élémentaires, la construction de $\widetilde{G}$ sur $G$ -- il nous
faut pour ceci ouvrir une parenthèse
\note{la page s'arrête ici, au premier tiers de la feuille ; le reste est blanc. C'est la trentième page depuis la page 2 : la mention « (29 p) » de la couverture de la page 1 couvre à une page près les pages 2 à 31, et la page 32 commence un autre texte.}

\page{33}
\note{la page 32 est une couverture rose, portant seulement, en haut à droite et au crayon, « (21 p) », d'une main qui n'est pas sûrement la sienne ; elle n'a pas de titre. Les pages 33 et 34 sont deux demi-feuilles, écrites dans le sens de la largeur.}
\[
  \begin{aligned}
    F_3 &= T \\
    F_4 &= T - 1 \\
    F_5 &= T^2 - T - 1 \\
    F_6 &= T - 2
  \end{aligned}
\]
On peut définir $F_n$ comme le polynôme minimal (sur $\mathbf{Q}$) de
$1 + 2\cos 2\pi/n$, \struck{\ill{}} ou encore de $1 + \zeta + \zeta^{-1}$, où
$\zeta$ est une racine \emph{primitive} $n$\add{-ième} de l'unité.
\note{ce premier bloc est à l'encre noire ; la liste commence à $F_3$. Un trait horizontal le sépare de ce qui suit, à l'encre bleue.}

\[
  \zeta_0 \xrightarrow[e_1]{} \zeta_1 \quad \zeta_2 \quad \zeta_3
\]
\[
  \zeta_1 - \zeta_0 = e_1 , \qquad \zeta_2 - \zeta_1 = e_2
\]
\[
  \begin{cases}
    u(\zeta_0) = \zeta_0 + e_1 \\
    u(e_1) = \zeta_2 - \zeta_1 = e_2 \\
    u(e_2) = \zeta_3 - \zeta_2 = -e_1 + (\alpha - 1)e_2
  \end{cases}
\]
\note{à droite, une figure : un cercle portant les points $\zeta_{-1}, \zeta_0, \zeta_1, \zeta_2, \zeta_3, \ldots, \zeta_{n-2}$ ; les vecteurs $e_1 = \zeta_1 - \zeta_0$, $e_2 = \zeta_2 - \zeta_1$, $e_3$ sont tracés comme côtés successifs d'un polygone régulier inscrit, et une flèche épaisse marquée $\alpha$ va de $\zeta_0$ à $\zeta_3$ . $u$ est la rotation qui envoie $\zeta_i$ sur $\zeta_{i+1}$. Le bas de la demi-feuille est déchiré.}

\page{34}
\[
  \mathrm{Alg\,Sp\acute{e}c}(\mathcal{M}) \simeq \mathbb{P}(\mathfrak{g}) \simeq \check{\mathbb{P}}(\widetilde{\mathfrak{g}}) \simeq \bigl[\check{P} = \mathrm{Dr}(P)\bigr] \simeq \mathrm{Div}_2^{+}(X)
\]
\note{sous « $\mathrm{Alg\,Sp\acute{e}c}(\mathcal{M})$ », un petit $\mathfrak{g}$ (?) ; au-dessus de $\mathbb{P}(\mathfrak{g})$ un tilde, qui fait peut-être lire $\mathbb{P}(\widetilde{\mathfrak{g}})$. Sous $\check{\mathbb{P}}(\widetilde{\mathfrak{g}})$, relié par $\wr$ : « Ss-gpes lisses commut.\ de dim 1 dans $G$ ». La feuille est déchirée en haut à droite.}

$\mathcal{T} \in \ldots(G)$
\note{la lettre ronde de sa main est rendue $\mathcal{T}$, comme au batch 3 ; ce qui la suit (« $\in$ », deux points, « $(G)$ ») se lit mal.}
\begin{itemize}
  \item[] $\underline{\mathrm{Cent}}(\mathcal{T}) = \mathcal{T}$ -- $\underline{\mathcal{T}}$ \uncertain{le} \uncertain{centralisateur} dans \ldots\ \marginal{\uncertain{(à tester)} des \ill{} de $\mathcal{T}$ telles que $g^2 \neq 1$ \ldots\ \ill{} \ill{} \ill{} \ill{} $\mathcal{T}$}
  \item[] $U = X - Y$ torseur sous $\mathcal{T}$
  \item[] \emph{Dans} $\mathcal{T}$\ \ $N(\mathcal{T})/\mathcal{T} \subset \underline{\mathrm{Aut}}(\mathcal{T})$ \struck{il y a la section} \ldots\ $= -\mathrm{id}_{\mathcal{T}}$ est dans $N(\mathcal{T})/\mathcal{T}$ \ldots\ \struck{\ill{}}
  \item[] \struck{les} \add{\uncertain{Sont}} images inverses \uncertain{\ldots} de $\overline{\mathcal{T}}$ dans $N(\mathcal{T})$ \struck{\ill{}} \ldots\ $\mathcal{T}$-torseur $X - Y$ \uncertain{est} \uncertain{comme} $X - Y$.
\end{itemize}
\note{les quatre lignes sont tenues à gauche par une accolade ; dans la marge gauche, en oblique, une note illisible où se lisent « \ldots\ $\mathcal{T}$ \ldots\ » et « \ldots\ gp.\ \ldots\ ». La colonne de droite est une note serrée qui court le long de ces lignes.}

\struck{\ill{} \ill{} $X - Y$, on \uncertain{a} \ill{} : \ill{}} \struck{\ill{}} \uncertain{qui}
\[
  (X - Y) \longrightarrow \overline{\mathcal{T}}
\]
\[
  \mathcal{T} \xrightarrow{\ 2\,\mathrm{id}_{\mathcal{T}}\ } \mathcal{T}
\]
\note{les deux avant-dernières lignes sont biffées ; la barre verticale dans $(X - Y) \to$ est un trait de séparation.}

\page{36}
\note{nouvelle numérotation des formules, qui repart de (1) ; c'est le début de la démonstration que les pages 41 à 44 (batch 3) achèvent.}
Soient $S$ schéma, \struck{$u \in M_2(S) = M_2(\Gamma(\mathcal{O}_S))$
\uncertain{supposons}} $V$ \struck{fibré} \add{Module} vectoriel de rang 2,
$u \in \underline{\mathrm{End}}(V)$
\[
  (1) \qquad \forall s \in S , \quad u_s \notin k(s) \cdot 1
\]
\note{la ligne « $V$ Module vectoriel de rang 2, $u \in \underline{\mathrm{End}}(V)$ » est écrite au-dessus de la ligne biffée ; « Module » remplace « fibré ».}
Considérons le commutant de $u$ dans $\underline{\mathrm{End}}(V)$ (le
foncteur). Je dis qu'il est égal au sous-fibré vectoriel
$\mathcal{O}_S\,\mathrm{id} + \mathcal{O}_S u$ de $\underline{\mathrm{End}}(V)$
-- donc le commutant dans $\mathrm{GL}(V)$ est formé des
$\alpha\,\mathrm{id} + \beta u$, avec
\[
  (2) \qquad \det(\alpha\,\mathrm{id} + \beta u) = \alpha^2 + \alpha\beta\, \mathrm{Tr}\, u + \beta^2 \det u \quad \text{inv.}
\]
\note{dans (2), devant $\alpha\beta$, un signe surchargé ; après $\det u$, un signe biffé.}
donc le commutant (ou stabilisateur) dans $\underline{\mathrm{GP}}(V)$
s'identifie au
\[
  C_{\underline{\mathrm{GP}}(V)}(u) \simeq \mathbb{P}^1 - Q_{\mathrm{Tr}\,u,\, \det u}
\]
où $Q_{\tau, \delta}$ désigne, de façon générale, la « quadrique » du
$\mathbb{P}^1$ définie par l'équation quadratique en $\alpha, \beta$
\[
  \alpha^2 + \tau\,\alpha\beta + \delta\,\beta^2
\]
\note{après $\beta^2$, un mot biffé ; dans l'indice de $Q_{\tau,\delta}$, $\tau$ est écrit sur un premier signe.}
\uncertain{Elle} \ill{} \ill{}, \ill{} un diviseur de degré 2 \struck{\ill{}}
qui est étale ssi
\[
  \tau^2 - 4\delta \quad \text{inv.}
\]
ce qui, dans le cas où $\tau = \mathrm{Tr}\,u$, $\delta = \det u$, signifie
que les val.\ propres de $u$ \uncertain{fibre} par fibre sont distinctes,
\note{le début de cette phrase ne se lit pas ; le sens attendu est que $Q_{\tau,\delta}$ est un diviseur de degré 2 de $\mathbb{P}^1$.}

\page{37}
Ce sous-groupe \add{$H \subset \underline{\mathrm{GP}}(V)$} (\uncertain{lisse}
\ill{}) ne dépend que du sous-module $\mathcal{O}_S + \mathcal{O}_S u$ de
$\underline{\mathrm{End}}(V)$, ou ce qui revient au même, de son image dans
$\underline{\mathrm{End}}(V)/\mathcal{O}_S \cdot 1 = \underline{\mathrm{Lie}}\,\underline{\mathrm{GP}}(V)$,
laquelle image est un sous-fibré inversible qui n'est autre que
$\mathrm{Lie}\, H$. Ceci précise relation entre sous-fibrés vectoriels
\add{$\mathfrak{h}$} de rang 1 de $\mathfrak{g} = \underline{\mathrm{GP}}(V) = G$
et sous-groupes \add{\uncertain{commutatifs}} \add{$H$} lisses de \struck{rang}
\add{dim rel.} 1 de $G$.
\note{les insertions ($H \subset \underline{\mathrm{GP}}(V)$, $\mathfrak{h}$, $H$, « dim rel. ») sont interlinéaires ; « commutatifs » est souligné d'un trait ondulé. L'égalité « $\mathfrak{g} = \underline{\mathrm{GP}}(V)$ » est écrite telle quelle, là où l'on attend $\mathrm{Lie}$.}

Est-il vrai que $H$ est le centralisateur de $\mathfrak{h}$
\add{$= \mathrm{Lie}\,\mathfrak{h}$} \uncertain{et} \uncertain{que} \ill{} de
$H$ \uncertain{forme} \uncertain{une} section \uncertain{unité}, \ill{} à
fortiori de $H$ lui-même ?
\note{« $= \mathrm{Lie}\,\mathfrak{h}$ » est inséré au-dessus de $\mathfrak{h}$, par une flèche ; la ligne se lit mal.}
\[
  g h g^{-1} = \lambda h \qquad \text{OPS} \quad \det u = -1
\]
\note{sous « $\det u = -1$ », un mot souligné d'un trait ondulé, lu avec doute « unipotent ».}
$\det h = \lambda^2 \det h$, \struck{$\mathrm{Tr}\,h = \lambda\,\mathrm{Tr}\,h$} \quad $\lambda^2 = 1$
d'où $\lambda = \pm 1$ \uncertain{loc.}\ \add{\uncertain{alors}}
\struck{si car.\ $\neq 2$, \ldots\ le cas $\pm 1$ \uncertain{disjoint} \ldots\ \uncertain{vient} le cas
$gug^{-1} = -u$, il est \ill{} possible que $u$ \ill{} \uncertain{soit} \ill{} (avec $\det u = -1$) soit \ill{} de la forme
$\mathrm{diag}(1, -1)$ \ill{} \ill{} \ill{} \ill{} \ldots\ de
\uncertain{valeurs} distinctes, \ill{} \ill{} \ill{} \ill{} \ldots\ est conjugué à \ill{}
réflexions dans $G$} donc $u^2 = 1$ \ill{} \ill{} \ill{} car.\ \uncertain{rés.}~2,
\note{à partir de « si car.\ $\neq 2$ », presque tout le bas de la page est annulé par des traits obliques et encadré ; à droite, en interligne, « $\mathrm{Tr}\,u = 0$, donc » et une matrice diagonale ; la matrice rendue $\mathrm{diag}(1, -1)$ est écrite en tableau. Seule la dernière ligne, « donc $u^2 = 1$ \ldots », est hors de la biffure.}

\page{38}
\note{ce début ne continue pas la page 37, dont la suite est la page 39 ; il semble prolonger la fin de la page 41 (batch 3), où $N/N^{\circ} \subset (\mathbf{Z}/p\mathbf{Z})^{\times}_S$ reste en suspens. Rattachement proposé par l'éditeur, non marqué par l'auteur.}
alors un s.-groupe étale \struck{qui \ill{} \ill{} fibres est \ill{} dans le
centre de la \ill{} des fibres} -- \add{\ill{}} il s'identifie à un sous-groupe
ouvert \struck{et fermé} de $(\mathbf{Z}/p\mathbf{Z})^{*}_S$, qui coïncide
avec le sous-groupe ouvert $\lbrace \pm 1\rbrace$ fibre par fibre, donc
\add{\uncertain{lui}} est égal.
\note{au-dessus de la première ligne, une insertion biffée : « et \ldots{} \ldots, donc \ldots ».}

La condition de normalisation, \add{sur une section $g \in G(S)$,} grâce au
régularité par changement de base, revient à la suivante : soit
$gug^{-1} = u$ (centralisateur), soit $gug^{-1} = u^{-1}$
\note{« sur une section $g \in G(S)$ » est inséré au-dessus de la ligne ; « $G(S)$ » se lit mal. Sous $u^{-1}$, un trait ondulé.}
$u = \begin{pmatrix} a & b \\ c & d \end{pmatrix}$ \qquad \struck{$ad - bc$ \uncertain{inv.}} \quad $\mathrm{GL}(2) \subset M_2$
\note{devant $\mathrm{GL}(2)$, un mot biffé ; $\subset$ est écrit verticalement, $M_2$ sous $\mathrm{GL}(2)$.}
\[
  v = \begin{pmatrix} x & y \\ z & t \end{pmatrix} \qquad
  v \longmapsto [\underline{u}, v] \qquad M_2 \xrightarrow{\ \mathrm{fu}\ } M_2
\]
\note{devant $v$, un mot biffé ; les lettres de la seconde ligne de $v$ sont lues avec doute. L'étiquette de la flèche, « fu », peut désigner $f_u$.}
\struck{$\mathcal{C}(u)$} est de rang $\leq 2$ [\uncertain{mineurs} d'ordre 3
\uncertain{tous} \uncertain{nuls}]

ssi \struck{\ill{}} est proj.\ de rang $= 2$ (\uncertain{mineurs} d'ordre 2)
\note{devant ces deux lignes, un mot biffé ; au début de la seconde, un symbole surchargé.}

\struck{\ill{}} $u$ \emph{pas mult.\ scalaire}
$\alpha + \beta u =$ \struck{$\alpha(\mathrm{id} + \alpha^{-1}\beta u)$}
\[
  \det(\alpha + \beta u) = \prod (\alpha + \beta\mu_i) = \alpha^2 + \alpha\beta\,\mathrm{Tr}\,u + \beta^2 \det u
\]
\note{les $\mu_i$ sont les valeurs propres de $u$ ; la lettre est mal formée. Le reste de la page est blanc.}

\page{39}
\[
  g h^2 g^{-1} = (\lambda h)^2 = \lambda^2 h^2 = h^2
\]
donc $g \in \mathcal{C}(h^2)$ donc OK.\ si $h^2 \neq 1$ i.e.\ $h$ non une
réflexion.
\note{le calcul prolonge celui de la page 37 ($ghg^{-1} = \lambda h$, $\lambda^2 = 1$) : la page 39 fait suite à la page 37, la page 38 s'intercalant hors de l'ordre de l'argument.}

Mais cela ne fait une belle jambe dans le cas (typique de car.~2 !) où
\struck{\ill{}} $H$ lui-même est tel que $h^2 = 1$ $\forall h \in \Gamma(U, H)$ \ldots !
Mais alors est-il possible d'avoir
\[
  g u g^{-1} = \lambda u^{-1} \qquad \forall \ \ldots
\]
pour $g$ fixé,
et $u$ \add{\uncertain{décrivant}} \struck{\ill{} fibré} formes $\alpha + \beta u_0$.
Prenant la section « universelle » $\alpha + \beta u_0$, où donc
$\alpha, \beta$ \uncertain{indéterminées} \ldots
\note{dans la formule, l'exposant $-1$ est surchargé ; « $\forall$ pour $g$ fixé » est écrit tel quel. « décrivant » est inséré au-dessus de la ligne.}
\[
  \underbrace{g(\alpha + \beta u_0)g^{-1}}_{\alpha + \beta\, g u_0 g^{-1} = \alpha + \beta\lambda_0 u_0} = \lambda(\alpha, \beta)\,(\alpha + \beta u_0)
\]
ce qui implique, \struck{\ill{} \ill{}} \struck{$gu_0g^{-1} = \lambda_0 u_0$}

$\alpha\,\lambda(\alpha, \beta) \equiv 1$ identité en $\alpha, \beta$ dans
$\mathcal{O}_S[\alpha, \beta]_{\Delta(\alpha, \beta)}$
\note{sous l'indice $\Delta(\alpha, \beta)$, une flèche vers « localisation ».}
d'où $\lambda(\alpha, \beta) = 1$, ok.

Donc $H$ est \struck{son propre normalisateur} \add{\uncertain{est-il}
\uncertain{son} \uncertain{propre} centralisateur de \ill{} \ill{} \ill{}
\uncertain{égal à} 1)}. Est-il le centralisateur \emph{de} $\mathfrak{h}$ ?
\note{la ligne « Donc $H$ est son propre normalisateur » est biffée d'un trait ondulé ; l'insertion au-dessous se lit mal.}

\page{40}
\[
  \begin{aligned}
    &g u g^{-1} = u + \lambda \\
    &\mathrm{Tr}\,u = \mathrm{Tr}\,u + 2\lambda
  \end{aligned}
\]
Si 2 inv.\ cela force $\lambda = 0$ -- \uncertain{gagné}.

\struck{Sinon, \ill{} \ill{} il y a contre-exemple,
$\det u = \det(u + \lambda) = \lambda^2 + (\mathrm{Tr}\,u)\lambda + \det u$, donc $\lambda(\mathrm{Tr}\,u + \lambda) = 0$ \ldots\ p.ex.\ la matrice
diagonale $u = (\zeta, \zeta + 1)$, $\zeta^2 + \zeta + 1 = 0$ \ldots\ est
équivalente à la matrice $(\zeta + 1, \zeta)$ \ldots\ sans être égale ! Ici le
centralisateur de $\mathfrak{h}$ dans $G$ (car.~2) est lisse de dim relative 1 !}
\note{tout ce bloc est annulé par un encadrement et de grands traits croisés ; la ligne « $\lambda^2(\ldots\mathrm{Tr}\,u + \lambda) = 0$ » y est de plus biffée, et plusieurs mots entre les formules se lisent mal.}

Si $u$ est diagonalisable, soit $u = (\zeta, \zeta')$ dans la translation par
$\zeta + \zeta'$ le transforme en $(\zeta', \zeta)$, qui lui est
\emph{conjugué}.
\note{« $\zeta + \zeta'$ » : lecture par le sens ; en car.~2, $u + (\zeta + \zeta')$ échange les deux valeurs propres.}

? Mais le centralisateur \emph{semble} \uncertain{être} \uncertain{donc} $H$,

\emph{NB} \struck{\ill{}} Dans $\mathfrak{g}$, algèbre de Lie de
$\mathrm{GP}(1)$ \add{avec la base habituelle}
\struck{le centralisateur de}
\[
  [H, X] = X , \quad [H, Y] = -Y , \quad [X, Y] = 2H = 0 \ \text{(en car.~2)}
\]
le centralisateur de $X$ \struck{\ill{}} \struck{\ill{}} \struck{\ill{} de $X$} est formé des
$xX + hH + yY$ tels que $h = 0$, donc est de dim trop grande -- celui de $H$
est OK\ldots
\note{sous « $= 0$ » : « \uncertain{en car.} 2 », lu avec doute. La page s'arrête ici ; la page 41 (batch 3) commence au milieu d'une démonstration sur un sous-groupe $\mu$ d'ordre premier $p$ impair.}

\end{document}
